% Problem record op_b08ad9d4371ed0cb. This TeX file is derived from
% database/problems_json/op_b08ad9d4371ed0cb.json by scripts/sync-tex.mjs; edit the JSON record
% and rerun the script rather than editing this file.
\section{Eight-ququart AME state and the ((7,4,4)) ququart MDS code}
\label{sec:op_b08ad9d4371ed0cb}

\paragraph{Problem.}
Does an absolutely maximally entangled state of eight ququarts exist?  More
precisely, with $[8]=\{1,\ldots,8\}$, determine whether there is a unit vector
$|\psi\rangle\in(\mathbb{C}^{4})^{\otimes 8}$ such that
\begin{equation}
  \operatorname{Tr}_{S^{c}}\!\left(|\psi\rangle\!\langle\psi|\right)
  =\frac{I_{4^{4}}}{4^{4}}
  \qquad\text{for every }S\subseteq[8]\text{ with }|S|=4.
  \label{eq:p40-ame-condition}
\end{equation}
Condition~\eqref{eq:p40-ame-condition} is equivalently the existence of a
rank-eight perfect tensor of bond dimension $4$: every flattening across a
$4|4$ partition is proportional to a unitary
\sourcecite{ref:p40-pastawski}{PYHP15}.  Under the pure-code correspondence,
it is also equivalent to a pure quantum MDS code with parameters
$\bigl[\!\bigl[8,0,5\bigr]\!\bigr]_{4}$
\sourcecite{ref:p40-shi}{SZZL26}.  Equivalently, the question asks whether
there is a quantum maximum-distance-separable code with parameters
$((7,4,4))_{4}$: a four-dimensional subspace
$\mathcal{Q}\subseteq(\mathbb{C}^{4})^{\otimes 7}$ whose orthogonal projector
$P$ satisfies
\begin{equation}
  PEP=c(E)\,P
  \quad\text{for every operator }E\text{ acting nontrivially on at most three ququarts},
  \qquad
  \dim\mathcal{Q}=4=4^{\,7-2(4-1)},
  \label{eq:p40-mds-code}
\end{equation}
where $c(E)\in\mathbb{C}$ and the last equality is saturation of the quantum
Singleton bound \sourcecite{ref:p40-huber-grassl}{HG20}.

\subsection*{Status}

Unsolved

\subsection*{Source}

Shi, Zhang, Zhao, and Li complete the seven-party AME classification and,
in their Conclusion and Discussion, explicitly identify
$\operatorname{AME}(8,4)$ as the first remaining homogeneous existence case
\sourcecite{ref:p40-shi}{SZZL26}.

\subsection*{Progress}

\begin{itemize}
  \item Bernal proved that a minimally supported $\operatorname{AME}(N,4)$
  state exists only for $N\leq 6$ (Theorem 6).  Thus any solution of
  \eqref{eq:p40-ame-condition} must have strictly more than $4^{4}$ nonzero
  coefficients in every local product basis
  \sourcecite{ref:p40-bernal}{Ber19}.

  \item Huber and Grassl proved that every quantum MDS code is pure
  (Theorem 5) and that a pure $((n,1,n/2+1))_{D}$ code, that is, an
  $\operatorname{AME}(n,D)$ state, exists exactly when an
  $((n-1,D,n/2))_{D}$ code exists (Observation 1 and Proposition 7).  For
  $n=8$ and $D=4$ this makes \eqref{eq:p40-ame-condition} and
  \eqref{eq:p40-mds-code} equivalent.  Given an orthonormal basis
  $\{|a_L\rangle\}_{a=0}^{3}$ of a code satisfying \eqref{eq:p40-mds-code},
  the state $|\psi\rangle=\tfrac12\sum_{a=0}^{3}|a\rangle_R|a_L\rangle$, with
  $R$ a reference ququart, satisfies \eqref{eq:p40-ame-condition};
  conversely, a Schmidt decomposition of such a state across one ququart
  versus the other seven yields the code.  Their
  shadow-inequality bounds hold for general, not necessarily stabilizer,
  quantum MDS codes but do not exclude
  $\bigl[\!\bigl[8,0,5\bigr]\!\bigr]_{4}$: in their table for local
  dimension $4$, the family with $n+k=8$ has upper bound
  $\bigl[\!\bigl[8,0,5\bigr]\!\bigr]_{4}$, while the highest distance
  then attained by a known code in that family is that of
  $\bigl[\!\bigl[6,2,3\bigr]\!\bigr]_{4}$ (Corollary 11 and Table 3)
  \sourcecite{ref:p40-huber-grassl}{HG20}.

  \item Ball, Moreno, and Simoens proved, with a computer-assisted
  classification, that no $\bigl[\!\bigl[7,1,4\bigr]\!\bigr]_{4}$
  stabilizer code exists, and deduced by projection that no
  $\bigl[\!\bigl[8,0,5\bigr]\!\bigr]_{4}$ stabilizer code exists
  (Section 5.2 and Lemma 49 of arXiv version 2).  Their stabilizers are
  abelian groups of Pauli operators labelled by $\mathbb{F}_{4}$; they
  correspond to additive codes in $\mathbb{F}_{4}^{2n}$ that are
  self-orthogonal under the trace-symplectic form and need not be
  $\mathbb{F}_{4}$-linear (Theorem 17), equivalently to binary stabilizer
  codes on $2n$ qubits grouped in pairs (Theorem 21).  Hence no state
  locally unitarily equivalent to an $\mathbb{F}_{4}$ stabilizer state
  satisfies \eqref{eq:p40-ame-condition}, and every code satisfying
  \eqref{eq:p40-mds-code} lies outside the $\mathbb{F}_{4}$ stabilizer
  formalism.  At the start of Section 5 they
  also remark that stabilizer codes with these parameters over the Pauli
  group labelled by $\mathbb{Z}/4\mathbb{Z}$ would yield binary
  $\bigl[\!\bigl[7,1,4\bigr]\!\bigr]_{2}$ and
  $\bigl[\!\bigl[8,0,5\bigr]\!\bigr]_{2}$ codes, which do not exist.
  These arguments do not address non-stabilizer states
  \sourcecite{ref:p40-ball}{BMS25}.

  \item W\'ojcik, Makuta, Bruzda, and Augusiak proved that no canonical
  $\mathbb{Z}_{d}$ graph state can be $\operatorname{AME}(N,d)$ when $4$
  divides $N$ and $d$ is even (Theorem 1).  This excludes canonical
  $\mathbb{Z}_{4}$ graph-state constructions for
  \eqref{eq:p40-ame-condition}.  Their argument uses the ring
  $\mathbb{Z}_{d}$ and does not apply to stabilizer states over
  $\mathbb{F}_{4}$: the authors point out that an $\mathbb{F}_{4}$
  stabilizer $\operatorname{AME}(4,4)$ state exists, although their theorem
  excludes a canonical $\mathbb{Z}_{4}$ graph state with those parameters
  \sourcecite{ref:p40-wojcik}{WMB+26}.

  \item Shi, Zhang, Zhao, and Li proved that
  $\operatorname{AME}(7,d)$ exists exactly for $d\geq 3$ (Theorem 1),
  settling the previously open cases $\operatorname{AME}(7,6)$ and
  $\operatorname{AME}(7,10)$ and, with earlier results, completing the
  homogeneous existence classification through seven parties.  They
  identify $\operatorname{AME}(8,4)$ as the first remaining open case
  (Conclusion and Discussion) \sourcecite{ref:p40-shi}{SZZL26}.
\end{itemize}

\subsection*{References}

\begin{enumerate}[label={},leftmargin=4.8em,labelsep=0.5em]
  \footnotesize
  \item[\textup{[PYHP15]}]\label{ref:p40-pastawski}
  F. Pastawski, B. Yoshida, D. Harlow, and J. Preskill,
  ``Holographic Quantum Error-Correcting Codes: Toy Models for the
  Bulk/Boundary Correspondence,'' \emph{Journal of High Energy Physics}
  \textbf{2015}(6), 149 (2015).
  \href{https://doi.org/10.1007/JHEP06(2015)149}{doi:10.1007/JHEP06(2015)149};
  \href{https://arxiv.org/abs/1503.06237}{arXiv:1503.06237}.

  \item[\textup{[Ber19]}]\label{ref:p40-bernal}
  A. Bernal,
  ``On the Existence of Absolutely Maximally Entangled States of Minimal
  Support II,'' \emph{Quantum Physics Letters} \textbf{8}, 1--4 (2019).
  \href{https://doi.org/10.18576/qpl/080101}{doi:10.18576/qpl/080101};
  \href{https://arxiv.org/abs/1807.00218}{arXiv:1807.00218}.

  \item[\textup{[HG20]}]\label{ref:p40-huber-grassl}
  F. Huber and M. Grassl,
  ``Quantum Codes of Maximal Distance and Highly Entangled Subspaces,''
  \emph{Quantum} \textbf{4}, 284 (2020).
  \href{https://doi.org/10.22331/q-2020-06-18-284}{doi:10.22331/q-2020-06-18-284};
  \href{https://arxiv.org/abs/1907.07733}{arXiv:1907.07733}.

  \item[\textup{[BMS25]}]\label{ref:p40-ball}
  S. Ball, E. Moreno, and R. Simoens,
  ``Stabilizer Codes Over Fields of Even Order,'' \emph{IEEE Transactions on
  Information Theory} \textbf{71}(5), 3707--3718 (2025); arXiv title
  ``Stabiliser codes over fields of even order.''
  \href{https://doi.org/10.1109/TIT.2024.3454480}{doi:10.1109/TIT.2024.3454480};
  \href{https://arxiv.org/abs/2401.06618}{arXiv:2401.06618}.

  \item[\textup{[WMB+26]}]\label{ref:p40-wojcik}
  J. W\'ojcik, O. Makuta, W. Bruzda, and R. Augusiak,
  ``On Non-Existence of Stabilizer Absolutely Maximally Entangled States in
  Even Local Dimensions,'' arXiv:2603.18193 (2026).
  \href{https://doi.org/10.48550/arXiv.2603.18193}{doi:10.48550/arXiv.2603.18193};
  \href{https://arxiv.org/abs/2603.18193}{arXiv:2603.18193}.

  \item[\textup{[SZZL26]}]\label{ref:p40-shi}
  F. Shi, X. Zhang, Q. Zhao, and L. Li,
  ``Complete Existence Classification of Seven-Partite Absolutely Maximally
  Entangled States,'' arXiv:2608.01011 (2026).
  \href{https://doi.org/10.48550/arXiv.2608.01011}{doi:10.48550/arXiv.2608.01011};
  \href{https://arxiv.org/abs/2608.01011}{arXiv:2608.01011}.
\end{enumerate}

\subsection*{Comment}

The unresolved question concerns arbitrary complex states satisfying
\eqref{eq:p40-ame-condition}, equivalently arbitrary, possibly non-additive,
codes satisfying \eqref{eq:p40-mds-code}.  Existing no-go results exclude minimal-support
states, all stabilizer states over $\mathbb{F}_{4}$ (including those given by
additive codes that are not $\mathbb{F}_{4}$-linear), and canonical
$\mathbb{Z}_{4}$ graph states; Ball, Moreno, and Simoens also remark that
stabilizer codes over $\mathbb{Z}/4\mathbb{Z}$ with these parameters do not
exist.  A construction must therefore have more than $4^{4}$ nonzero
coefficients in every local product basis and lie, up to local unitaries,
outside both stabilizer formalisms, whereas a nonexistence proof must treat
general quantum codes, for which the shadow-inequality bounds of Huber and
Grassl do not suffice.  A literature check through 15 September 2026 found
no such construction or proof.

\subsection*{Field}

Quantum Resource Theory; Quantum Error Correction

\subsection*{Topic}

Absolutely maximally entangled states; Quantum coding theory

\subsection*{ID}

\texttt{op\_b08ad9d4371ed0cb}
